% Morino Dirichlet formulation: surface S, wake W, interior potential zero, Kutta condition.
\begin{tikzpicture}[font=\small]
  \def\xe{8}
  \path[fill=cfill] (0,0) .. controls (0,1.7) and (4,2.1) .. (\xe,0) .. controls (4,-2.1) and (0,-1.7) .. (0,0);
  \draw[cblue, very thick] (0,0) .. controls (0,1.7) and (4,2.1) .. (\xe,0) .. controls (4,-2.1) and (0,-1.7) .. (0,0);
  \node[cink, align=center] at (3.8,0) {内部 $V_i$\\$\Phi_i=0$};
  % wake
  \draw[corange, very thick, dashed] (\xe,0) -- (11.6,0);
  \draw[corange, very thick] (\xe,0.04) -- (11.6,0.04);
  \node[corange, align=center] at (10.0,-0.75) {後流 $W$\\{\footnotesize 二重層のみ}\\$\mu_w=\mu_u-\mu_l$};
  % free stream
  \foreach \y in {-2.4,-1.2,0,1.2,2.4}{
    \draw[cink2, ->] (-3.2,\y) -- (-2.2,\y);
  }
  \node[cink2] at (-2.7,3.1) {$\bm U$（一様流）};
  % normal at a surface point
  \coordinate (P) at (4.0,1.52);
  \draw[cink, ->, thick] (P) -- ++(0.15,0.85) node[above right] {$\bm n$};
  \fill[cink] (P) circle (1.6pt);
  \node[cblue, align=left, anchor=west] at (3.0,3.5) {$S$ 上：doublet $\mu$、source $\sigma=-\bm n\!\cdot\!\bm U$\\{\footnotesize 外側の $\Phi=\mu$（内部 $\Phi_i=0$ との差）}};
  \node[cblue] at (-0.55,1.7) {表面 $S$};
  \draw[cblue, ->] (-0.4,1.5) -- (0.25,0.9);
  % trailing edge panels
  \fill[cblue] (\xe-0.45,0.17) circle (1.6pt) node[above, font=\footnotesize] {$\mu_u$};
  \fill[cblue] (\xe-0.45,-0.17) circle (1.6pt) node[below, font=\footnotesize] {$\mu_l$};
  \node[cink2, align=center, font=\footnotesize] at (\xe+0.6,1.5) {後縁\\Kutta 条件};
  \draw[cink2, ->] (\xe+0.5,1.0) -- (\xe,0.12);
  \node[cink2] at (5.3,-3.0) {外部 $V_e$（流体領域）};
  \node[cink2, font=\footnotesize] at (11.0,-2.9) {遠方で $\Phi\to0$};
\end{tikzpicture}
